\section{The layer lemma and the reduction to exactly \texorpdfstring{$mt$}{mt} points}\label{ll:sec}

This section gives written proofs of the declarations of the Lean file \dirn{LAYERLEAN/LayerLemma.lean} (namespace
\lean{ErdosSzekeres.LayerLemma}) that concern hull vertices, orbits of a map of finite period, the layer lemma, the
peeling of convex layers, and the reduction of the sixfold and sevenfold octagon bounds to statements about sets of
exactly $mt$ points. The main results are:
\begin{itemize}
\item the layer lemma in abstract form (Theorem~\ref{ll:thm:core}): if an injective affine map $f$ sends a finite set
  $S$ into itself, every point of $S$ has minimal period $m>0$ under $f$, and no $K$ points of $S$ are in convex
  position for some $K\le 2m$, then for every hull vertex $x$ of $S$ the hull vertices of $S$ are exactly the points
  $x,f(x),\dots,f^{m-1}(x)$, so that, if $S$ is nonempty, they form a single $f$-orbit; no general position is
  assumed;
\item its rotation form (Theorem~\ref{ll:thm:core-rot}), the description of all convex layers
  (Theorem~\ref{ll:thm:layers}, Lemma~\ref{ll:lem:exists-layer}) and of their sizes (Theorem~\ref{ll:thm:card-peel});
\item the reduction to exactly $mt$ points (Theorem~\ref{ll:thm:reduction}) and the sixfold and sevenfold octagon
  theorems (Theorems~\ref{ll:thm:sixfold} and~\ref{ll:thm:sevenfold}). The last two are \emph{conditional}: they take
  the propositions \lean{Sixfold66} and \lean{Sevenfold70} (Definition~\ref{ll:def:66-70}) as hypotheses, and these
  propositions are not proved in Lean (Remark~\ref{ll:rem:hyp}).
\end{itemize}

\paragraph{Sources.} The Lean source is \dirn{LAYERLEAN/LayerLemma.lean} (1148 lines, SHA-256
\texttt{afb8542dceb0}\dots). This section treats its lines 1--469 and 1034--1148; the geometric lemmas of lines
471--1032 (strict nesting, the centre for even $m$, chains, polar rules, radius order) belong to other sections (for
example Theorem~\ref{lg:thm:nesting} and Lemma~\ref{lg:lem:centre}). The notions of point, general position, convex
position and \lean{HasConvexNGon} come from \dirn{ErdosSzekeres/Definitions.lean} (SHA-256 \texttt{2425500eefd1}\dots)
and are recalled in Definitions~\ref{fd:def:point}, \ref{fd:def:gp}, \ref{fd:def:cp} and~\ref{fd:def:ngon}.

\paragraph{Lean evidence.} The file was checked with Lean 4.24.0 and Mathlib \texttt{f897ebcf} (the repository's pins)
by \texttt{lake env lean}, run from the repository root. It is not a module of the default build target, so a plain
\texttt{lake build} does not re-check it. Two records show that the current file compiles. The review report
\dirn{LEANEXT/review/REPORT.md} of 4 November 2026 (item~1) records exit code~0 and no messages for the reviewed file,
which it identifies by its SHA-256 \texttt{afb8542d}\dots{} and its 1148 lines. The re-check
\dirn{LAYERLEAN/logs/axioms_recheck_20261106/out.txt} of 6 November 2026 records exit status~0 for a file whose first
1148 lines are byte-identical to \dirn{LAYERLEAN/LayerLemma.lean}, followed by 19 \texttt{\#print axioms} commands;
its only Lean output is the 19 printed axiom lines. The log \dirn{LAYERLEAN/logs/build.log} (exit code~0, no
messages) names no file hash; it was committed with the first version of the file (454 lines, SHA-256
\texttt{804b1d25f79a}\dots, commit \texttt{6429e9a70}) and was not updated by the five later commits that changed the
file, so we do not use it. Axioms: \dirn{LAYERLEAN/logs/axioms.log} (28 theorems) and the re-check output just named
(19 theorems) report exactly the axioms \lean{propext}, \lean{Classical.choice} and \lean{Quot.sound} for every
printed theorem. The review log
\dirn{LEANEXT/review/axioms.log} prints 59 theorems, namely all theorems of the file that are not \texttt{@[simp]}
lemmas, among them every theorem of this section except the seven \texttt{@[simp]} lemmas of
Lemma~\ref{ll:lem:complex}(a)--(b); it reports the same three axioms for each printed theorem, except that
\lean{iterate_mem} depends only on \lean{propext} and \lean{Quot.sound}. This log comes from the review whose report is
named above; the check file from which the axioms were printed is not kept in the record, so the identification of
the printed theorems with those of this file rests on the report. The seven \texttt{@[simp]} lemmas are printed in
none of the three logs; their proofs (lines 236--255) are
\texttt{rfl}, extensionality (\lean{Complex.ext}, \texttt{funext}) and, for \lean{toC_crot}, the lemma
\lean{toC_ofC}. The file contains no \texttt{sorry}, \texttt{admit}, \texttt{axiom} or \texttt{native\_decide}, and its
only \texttt{decide} (line 1125) proves that $6$ is even and is evaluated by the kernel. So no statement of this
section depends on \lean{sorryAx}, on \lean{Lean.ofReduceBool} (\texttt{native\_decide}) or on an axiom of the
project. Every lemma and theorem below whose head names a Lean declaration is \emph{Lean-checked} in this sense
(Theorems~\ref{ll:thm:sixfold} and~\ref{ll:thm:sevenfold} as implications from their hypotheses); for
Lemma~\ref{ll:lem:rot}(b) the Lean declaration is the definition \lean{rotAff}, whose affine-map property is a proof
obligation discharged inside the definition and checked with it. The definitions whose heads name a Lean declaration restate Lean
definitions, which the same build type-checks; being definitions, they are not Lean-checked claims. In particular,
Definition~\ref{ll:def:66-70} only introduces the propositions \lean{Sixfold66} and \lean{Sevenfold70}, which are not
proved (Remark~\ref{ll:rem:hyp}). The proofs below are \emph{written proofs (not refereed)}; for each statement that
corresponds to a Lean declaration we say whether its proof follows the Lean proof.

\paragraph{Dependencies.} From the foundations section we use Definitions~\ref{fd:def:point}, \ref{fd:def:orient},
\ref{fd:def:gp}, \ref{fd:def:cp}, \ref{fd:def:ngon}, \ref{fd:def:inv}, \ref{fd:def:rot}, Lemmas~\ref{fd:lem:cp-hull},
\ref{fd:lem:mono}, \ref{fd:lem:collinear} and~\ref{fd:lem:affine-cp}, and the lemma of that section titled ``Convex
hulls'' (in (C3) below). Theorem~\ref{ll:thm:sixfold} uses Lemma~\ref{lg:lem:centre} (exclusion of the centre for
even $m$); for the case needed we also include the short argument. Nothing else from other sections is used.

\subsection{Conventions}

Points are elements of $\R^2$ (Lean \lean{Point}, defined as \texttt{Fin 2} $\to\R$), written $p=(p_0,p_1)$. We use
the following conventions and standard facts.
\begin{enumerate}[label=(C\arabic*)]
\item \emph{Finite sets.} A finite set of points is a Lean \texttt{Finset Point}; $\card S$ is its cardinality;
  $S\setminus\{c\}$ is Lean \texttt{S.erase c}; ``$S$ is nonempty'' is Lean \lean{S.Nonempty}.
\item \emph{Iterates and minimal period.} For a set $\alpha$ and $f:\alpha\to\alpha$ put $f^0=\mathrm{id}$ and
  $f^{i+1}=f\circ f^i$. Lean's \texttt{f\textasciicircum[i]} (\lean{Nat.iterate}) is defined by the other recursion,
  $f^{i+1}=f^i\circ f$; an easy induction shows that both recursions define the same maps (Mathlib's
  \lean{Function.iterate_succ'}), and $f^{i+j}=f^i\circ f^j$ (Mathlib's \lean{Function.iterate_add}). The
  \emph{minimal period}
  $\operatorname{per}_f(x)$ of $x\in\alpha$ is Mathlib's \texttt{Function.minimalPeriod f x}: the least $n\ge1$ with
  $f^n(x)=x$ if such an $n$ exists, and $0$ otherwise. Directly from this definition:
  \begin{itemize}
  \item[(P1)] $f^{\operatorname{per}_f(x)}(x)=x$ for every $x$ (if no period exists, $f^0(x)=x$);
  \item[(P2)] if $n\ge1$ and $f^n(x)=x$, then $1\le\operatorname{per}_f(x)\le n$.
  \end{itemize}
\item \emph{Convex hulls.} For $A\subseteq\R^2$, $\conv(A)$ is the smallest convex set containing $A$ (Mathlib
  \lean{convexHull}). It is the set of all convex combinations $\sum_{i=1}^r\lambda_ia_i$ with $r\ge1$, $a_i\in A$,
  $\lambda_i\ge0$ and $\sum_i\lambda_i=1$. This is part~(b) of the lemma titled ``Convex hulls'' in the foundations
  section, where it is proved for the whole document (see also R.~T.~Rockafellar, \emph{Convex Analysis}, Princeton
  1970, Theorem~2.3). Hence: $A\subseteq\conv(A)$; $A\subseteq B$ implies $\conv(A)\subseteq\conv(B)$; if
  $A\subseteq C$ with $C$ convex
  then $\conv(A)\subseteq C$, so $\conv(\conv(A))=\conv(A)$; $\conv(\emptyset)=\emptyset$; a convex combination of
  points of a convex set $C$ lies in $C$. Merging equal points and deleting zero coefficients, every point of
  $\conv(A)$ is a convex combination of pairwise distinct points of $A$ with positive coefficients and at least one
  term.
\item \emph{Extreme points.} A point $x$ is an \emph{extreme point} of $A\subseteq\R^2$ (Mathlib
  \lean{Set.extremePoints}) if $x\in A$ and, whenever $x=ax_1+bx_2$ with $x_1,x_2\in A$, $a,b>0$ and $a+b=1$,
  we have $x_1=x$. This is Mathlib's definition at the pinned revision (with $x$ in the open segment between $x_1$ and
  $x_2$); exchanging the roles of $x_1$ and $x_2$ shows that then also $x_2=x$.
\item \emph{Affine maps.} A Lean affine map $f$, an element of the type \texttt{Point}~$\to^{\mathrm a}[\R]$~\texttt{Point},
  consists of a map $f:\R^2\to\R^2$ and a linear map $L$ with $f(v+p)=Lv+f(p)$ for all $v,p$; taking $p=0$,
$f(v)=Lv+b$ with $b=f(0)$, and conversely every map
  $p\mapsto Lp+b$ with $L$ linear is of this form. ``$f$ is injective'' refers to the underlying map on all of $\R^2$.
  If $\sum_i\lambda_i=1$ then $f(\sum_i\lambda_ia_i)=\sum_i\lambda_i(La_i+b)=\sum_i\lambda_if(a_i)$; hence, by (C3),
  $\conv(f(A))\subseteq f(\conv(A))$ for every $A\subseteq\R^2$ (this inclusion is part of
  Lemma~\ref{fd:lem:affine-cp}).
\item \emph{Convex position.} Convex position, general position and \lean{HasConvexNGon} are as in
  Definitions~\ref{fd:def:cp}, \ref{fd:def:gp} and~\ref{fd:def:ngon}: \lean{HasConvexNGon} $K$ $S$ means that some
  $T\subseteq S$ with $\card T=K$ is in convex position, and ``$S$ is $K$-free'' means $\neg\,$\lean{HasConvexNGon}
  $K$ $S$. We use Lemma~\ref{fd:lem:cp-hull} ($S$ is in convex position iff $p\notin\conv(S\setminus\{p\})$ for every
  $p\in S$) and Lemma~\ref{fd:lem:mono} (monotonicity). In particular the empty set is in convex position, so
  \lean{HasConvexNGon} $0$ $S$ holds for every $S$.
\item \emph{Angles.} In Lean the angle $2\pi/m$ is \texttt{2 * Real.pi / m} with $m\in\N$ cast to $\R$; division by
  $0$ gives $0$. For $t\in\N$ the subtraction $t\mathbin{\dot{-}}1$ is truncated ($0\mathbin{\dot{-}}1=0$).
\end{enumerate}

\subsection{Orbit segments}

\begin{definition}[orbit segment; \leanref{ErdosSzekeres.LayerLemma.orbitFS}{LAYERLEAN/LayerLemma.lean:48}]
\label{ll:def:orbit}
Let $\alpha$ be a set (Lean: a type with decidable equality), $f:\alpha\to\alpha$, $p\in\N$ and $x\in\alpha$. The
\emph{orbit segment} is the finite set
\[\operatorname{orb}_p(f,x)=\{f^i(x):\ i\in\N,\ 0\le i<p\}\]
(Lean \texttt{orbitFS f p x}, the image of $\{0,\dots,p-1\}$ under $i\mapsto f^i(x)$). It is empty for $p=0$.
\end{definition}

\begin{lemma}[orbit segments; \leanref{ErdosSzekeres.LayerLemma.iterate_mem}{LAYERLEAN/LayerLemma.lean:52},
\leanref{ErdosSzekeres.LayerLemma.mem_orbitFS_self}{LAYERLEAN/LayerLemma.lean:59},
\leanref{ErdosSzekeres.LayerLemma.orbitFS_subset}{LAYERLEAN/LayerLemma.lean:62},
\leanref{ErdosSzekeres.LayerLemma.orbitFS_invariant}{LAYERLEAN/LayerLemma.lean:68},
\leanref{ErdosSzekeres.LayerLemma.card_orbitFS_of_minimalPeriod}{LAYERLEAN/LayerLemma.lean:80},
\leanref{ErdosSzekeres.LayerLemma.mem_orbitFS_of_not_disjoint}{LAYERLEAN/LayerLemma.lean:91}]\label{ll:lem:orbit}
Let $\alpha$ be a set and $f:\alpha\to\alpha$.
\begin{enumerate}[label=(\alph*)]
\item (\lean{iterate_mem}) If $S\subseteq\alpha$ is finite, $f(y)\in S$ for every $y\in S$, and $x\in S$, then
  $f^i(x)\in S$ for every $i\in\N$.
\item (\lean{mem_orbitFS_self}) If $0<p$, then $x\in\operatorname{orb}_p(f,x)$ for every $x\in\alpha$.
\item (\lean{orbitFS_subset}) If $S\subseteq\alpha$ is finite, $f(y)\in S$ for every $y\in S$, and $x\in S$, then
  $\operatorname{orb}_p(f,x)\subseteq S$ for every $p\in\N$.
\item (\lean{orbitFS_invariant}) If $0<p$ and $f^p(x)=x$, then $f(y)\in\operatorname{orb}_p(f,x)$ for every
  $y\in\operatorname{orb}_p(f,x)$.
\item (\lean{card_orbitFS_of_minimalPeriod}) If $\operatorname{per}_f(x)=p$, then
  $\card{\operatorname{orb}_p(f,x)}=p$.
\item (\lean{mem_orbitFS_of_not_disjoint}) If $0<p$, $f^p(x)=x$, $f^p(y)=y$ and the sets $\operatorname{orb}_p(f,x)$
  and $\operatorname{orb}_p(f,y)$ are not disjoint, then $y\in\operatorname{orb}_p(f,x)$.
\end{enumerate}
\end{lemma}

\begin{proof}
We follow the Lean proofs.
(a) Induction on $i$: $f^0(x)=x\in S$, and if $f^i(x)\in S$ then $f^{i+1}(x)=f(f^i(x))\in S$.
(b) Take $i=0<p$: $f^0(x)=x$.
(c) Every element of $\operatorname{orb}_p(f,x)$ is some $f^i(x)$, which lies in $S$ by (a).
(d) Let $y=f^i(x)$ with $0\le i<p$; then $f(y)=f^{i+1}(x)$. If $i+1<p$ this lies in $\operatorname{orb}_p(f,x)$.
Otherwise $i+1=p$ and $f(y)=f^p(x)=x\in\operatorname{orb}_p(f,x)$ by (b).
(e) If $p=0$, then $\operatorname{orb}_0(f,x)=\emptyset$ has $0$ elements. Let $p\ge1$. It suffices to show that
$f^0(x),\dots,f^{p-1}(x)$ are pairwise distinct. By (P1), $f^p(x)=x$. Suppose $f^i(x)=f^j(x)$ with $0\le i<j\le p-1$.
Applying $f^{p-j}$ gives
\[f^{p-j+i}(x)=f^{p-j}(f^i(x))=f^{p-j}(f^j(x))=f^p(x)=x .\]
Since $j\le p-1$ and $i\le j-1$, we have $1\le p-j+i\le p-1$, so (P2) gives $\operatorname{per}_f(x)\le p-j+i<p$, a
contradiction.
(f) Let $z$ lie in both sets, say $z=f^j(y)$ with $0\le j<p$. Then $f^{p-j}(z)=f^{p}(y)=y$. By (d), $f$ maps the finite
set $\operatorname{orb}_p(f,x)$ into itself, and $z\in\operatorname{orb}_p(f,x)$; so by (a), applied to this set,
$y=f^{p-j}(z)\in\operatorname{orb}_p(f,x)$.
\end{proof}

\begin{definition}[invariance; \leanref{ErdosSzekeres.LayerLemma.Invariant}{LAYERLEAN/LayerLemma.lean:105}]
\label{ll:def:inv}
For a map $f:\R^2\to\R^2$ and a finite set $S\subseteq\R^2$, $S$ is \emph{$f$-invariant} (Lean
\texttt{Invariant f S}) if $f(p)\in S$ for every $p\in S$, that is, $f$ maps $S$ \emph{into} itself. This is the notion of
Definition~\ref{fd:def:inv}. (For injective $f$ it implies $f(S)=S$; see the proof of Lemma~\ref{ll:lem:map-hv}.)
\end{definition}

\subsection{Hull vertices}

\begin{definition}[hull vertices; \leanref{ErdosSzekeres.LayerLemma.hullVerts}{LAYERLEAN/LayerLemma.lean:111}]
\label{ll:def:hv}
For a finite set $S\subseteq\R^2$ let
\[\hv(S)=\{x\in S:\ x\notin\conv(S\setminus\{x\})\}\]
(Lean \texttt{hullVerts S}, the subset of $S$ cut out by this condition). Its elements are the \emph{hull vertices} of
$S$.
\end{definition}

\begin{lemma}[membership; \leanref{ErdosSzekeres.LayerLemma.mem_hullVerts}{LAYERLEAN/LayerLemma.lean:114},
\leanref{ErdosSzekeres.LayerLemma.hullVerts_subset}{LAYERLEAN/LayerLemma.lean:118}]\label{ll:lem:hv-basic}
Let $S$ be a finite set of points. (a) For every point $x$: $x\in\hv(S)$ if and only if $x\in S$ and
$x\notin\conv(S\setminus\{x\})$. (b) $\hv(S)\subseteq S$.
\end{lemma}

\begin{proof}
As in Lean: (a) is Definition~\ref{ll:def:hv}, and (b) follows from (a).
\end{proof}

\begin{lemma}[hull vertices; \leanref{ErdosSzekeres.LayerLemma.convexPosition_hullVerts}{LAYERLEAN/LayerLemma.lean:121},
\leanref{ErdosSzekeres.LayerLemma.extreme_mem_hullVerts}{LAYERLEAN/LayerLemma.lean:130},
\leanref{ErdosSzekeres.LayerLemma.subset_convexHull_hullVerts}{LAYERLEAN/LayerLemma.lean:146},
\leanref{ErdosSzekeres.LayerLemma.hullVerts_nonempty}{LAYERLEAN/LayerLemma.lean:139}]\label{ll:lem:hv}
Let $S$ be a finite set of points.
\begin{enumerate}[label=(\alph*)]
\item (\lean{convexPosition_hullVerts}) $\hv(S)$ is in convex position.
\item (\lean{extreme_mem_hullVerts}) Every extreme point of $\conv(S)$ belongs to $\hv(S)$.
\item (\lean{subset_convexHull_hullVerts}) $S\subseteq\conv(\hv(S))$.
\item (\lean{hullVerts_nonempty}) If $S$ is nonempty, then $\hv(S)$ is nonempty.
\end{enumerate}
\end{lemma}

\begin{proof}
The proof of (a) follows Lean. For (b)--(d) we give elementary proofs, different from the Lean proofs. Lean proves (b)
with two Mathlib facts: the extreme points of $\conv(S)$ lie in $S$, and ``$x$ is an extreme point of a convex set $C$
iff $x\in C$ and $C\setminus\{x\}$ is convex''. It deduces (d) from (b) and the existence of an extreme point of the
nonempty compact set $\conv(S)$ (Mathlib's Krein--Milman lemma), and (c) from (b) and the Krein--Milman theorem
($\conv(S)$ is the closure of the convex hull of its extreme points).

(a) By Lemma~\ref{fd:lem:cp-hull} we must show $x\notin\conv(\hv(S)\setminus\{x\})$ for every $x\in\hv(S)$. Since
$\hv(S)\setminus\{x\}\subseteq S\setminus\{x\}$, (C3) gives $\conv(\hv(S)\setminus\{x\})\subseteq\conv(S\setminus\{x\})$,
and $x\notin\conv(S\setminus\{x\})$ because $x\in\hv(S)$.

(b) Let $x$ be an extreme point of $C=\conv(S)$. We first show:
\begin{itemize}
\item[$(*)$] if $x=\sum_{i=1}^r\lambda_is_i$ with $r\ge1$, $s_i\in C$, $\lambda_i>0$ and $\sum_i\lambda_i=1$, then
  $s_i=x$ for every $i$.
\end{itemize}
If $r=1$ then $\lambda_1=1$ and $s_1=x$. Let $r\ge2$ and fix $j$. Then $0<\lambda_j<1$, since some other coefficient is
positive. The point $y_j=\sum_{i\ne j}\frac{\lambda_i}{1-\lambda_j}s_i$ is a convex combination of points of the convex
set $C$, so $y_j\in C$ by (C3), and $x=\lambda_js_j+(1-\lambda_j)y_j$. By (C4), applied with $x_1=s_j$, $x_2=y_j$,
$a=\lambda_j$, $b=1-\lambda_j$, we get $s_j=x$. This proves $(*)$.

Since $x\in C$, by (C3) $x=\sum_{i=1}^r\lambda_is_i$ with $r\ge1$, $s_i\in S\subseteq C$, $\lambda_i>0$,
$\sum_i\lambda_i=1$; by $(*)$, $x=s_1\in S$. Suppose $x\in\conv(S\setminus\{x\})$. Then $x=\sum_{i=1}^r\mu_it_i$ with
$r\ge1$, $t_i\in S\setminus\{x\}\subseteq C$, $\mu_i>0$, $\sum_i\mu_i=1$, and $(*)$ gives $t_1=x$, contradicting
$t_1\ne x$. Hence $x\notin\conv(S\setminus\{x\})$, and $x\in\hv(S)$.

(c) Strong induction on $\card S$. If $\hv(S)=S$, then $S\subseteq\conv(S)=\conv(\hv(S))$; this covers $S=\emptyset$.
Otherwise choose $y\in S\setminus\hv(S)$; then $y\in\conv(S')$ where $S'=S\setminus\{y\}$. As $\card{S'}<\card S$, the
induction hypothesis gives $S'\subseteq\conv(\hv(S'))$.

We claim $\hv(S')\subseteq\hv(S)$. Let $z\in\hv(S')$, so $z\in S'$ (in particular $z\in S$ and $z\ne y$) and
$z\notin\conv(A)$, where $A=S'\setminus\{z\}=S\setminus\{y,z\}$. Suppose $z\in\conv(S\setminus\{z\})$. Since
$S\setminus\{z\}=A\cup\{y\}$ and $S'=A\cup\{z\}$, (C3) and merging of equal points give
\[z=\sum_{a\in A}\alpha_aa+\lambda y,\qquad y=\sum_{a\in A}\beta_aa+\mu z,\]
with all $\alpha_a,\beta_a,\lambda,\mu\ge0$, $\sum_a\alpha_a+\lambda=1$ and $\sum_a\beta_a+\mu=1$ (the second
representation exists because $y\in\conv(S')$). Substituting the second into the first,
\[(1-\lambda\mu)\,z=\sum_{a\in A}(\alpha_a+\lambda\beta_a)\,a .\]
If $\lambda\mu<1$, the coefficients $(\alpha_a+\lambda\beta_a)/(1-\lambda\mu)$ are nonnegative with sum
$\bigl((1-\lambda)+\lambda(1-\mu)\bigr)/(1-\lambda\mu)=1$, so $z\in\conv(A)$, a contradiction. If $\lambda\mu=1$, then
$\lambda=\mu=1$ because $\lambda,\mu\in[0,1]$; so all $\alpha_a=0$ and $z=y$, contradicting $z\ne y$. This proves the
claim.

By the claim and (C3), $S'\subseteq\conv(\hv(S'))\subseteq\conv(\hv(S))$, and then
$y\in\conv(S')\subseteq\conv(\conv(\hv(S)))=\conv(\hv(S))$. Hence $S=S'\cup\{y\}\subseteq\conv(\hv(S))$.

(d) If $\hv(S)=\emptyset$, then by (c) $S\subseteq\conv(\emptyset)=\emptyset$, so $S$ is empty.
\end{proof}

\begin{lemma}[affine invariance of hull vertices;
\leanref{ErdosSzekeres.LayerLemma.map_mem_hullVerts}{LAYERLEAN/LayerLemma.lean:162}]\label{ll:lem:map-hv}
Let $f$ be an affine map of $\R^2$ (Lean: $f$ of type \texttt{Point}~$\to^{\mathrm a}[\R]$~\texttt{Point}) that is
injective, and let $S$ be a finite set
with $f(y)\in S$ for every $y\in S$. If $x\in\hv(S)$, then $f(x)\in\hv(S)$.
\end{lemma}

\begin{proof}
We follow the Lean proof. (1) $f(S)=S$: as $f$ is injective, $\card{f(S)}=\card S$, and $f(S)\subseteq S$ with $S$
finite. (2) $f(x)\in S$, since $x\in S$. (3) Suppose $f(x)\in\conv(S\setminus\{f(x)\})$. Every $y\in S$ with
$y\ne f(x)$ is, by (1), of the form $y=f(z)$ with $z\in S$, and $z\ne x$ since $y\ne f(x)$; thus
$S\setminus\{f(x)\}\subseteq f(S\setminus\{x\})$. By (C3) and (C5),
\[f(x)\in\conv\bigl(f(S\setminus\{x\})\bigr)\subseteq f\bigl(\conv(S\setminus\{x\})\bigr),\]
so $f(x)=f(u)$ for some $u\in\conv(S\setminus\{x\})$. Injectivity of $f$ on all of $\R^2$ (note that $u$ need not lie
in $S$) gives $x=u\in\conv(S\setminus\{x\})$, contradicting $x\in\hv(S)$. By Lemma~\ref{ll:lem:hv-basic}(a), (2) and
(3) give $f(x)\in\hv(S)$.
\end{proof}

\subsection{The core lemma}

\begin{theorem}[layer lemma, abstract form;
\leanref{ErdosSzekeres.LayerLemma.hullVerts_eq_orbitFS}{LAYERLEAN/LayerLemma.lean:189}]\label{ll:thm:core}
Let $f$ be an affine map of $\R^2$ (Lean: $f$ of type \texttt{Point}~$\to^{\mathrm a}[\R]$~\texttt{Point}) that is
injective, let $S$ be a finite set of
points, and let $m,K\in\N$ with $0<m$ and $K\le 2m$. Assume:
\begin{enumerate}[label=(\roman*)]
\item $f(y)\in S$ for every $y\in S$;
\item $\operatorname{per}_f(y)=m$ for every $y\in S$;
\item $S$ is $K$-free, i.e.\ $\neg\,$\lean{HasConvexNGon} $K$ $S$.
\end{enumerate}
Then for every $x\in\hv(S)$,
\[\hv(S)=\operatorname{orb}_m(f,x)=\{f^i(x):\ 0\le i<m\}.\]
\end{theorem}

\begin{proof}
We follow the Lean proof. Write $H=\hv(S)$, a finite set with $H\subseteq S$ (Lemma~\ref{ll:lem:hv-basic}).

Step 1. $f(y)\in H$ for every $y\in H$, by Lemma~\ref{ll:lem:map-hv} and (i).

Step 2. $f^m(y)=y$ for every $y\in S$, by (ii) and (P1).

Step 3. $\operatorname{orb}_m(f,x)\subseteq H$, by Lemma~\ref{ll:lem:orbit}(c) applied to the finite set $H$
(Step~1) and $x\in H$.

Step 4. $H\subseteq\operatorname{orb}_m(f,x)$. Let $y\in H$ and suppose $y\notin\operatorname{orb}_m(f,x)$. Both $x$ and
$y$ lie in $S$, so $f^m(x)=x$ and $f^m(y)=y$ by Step~2. If $\operatorname{orb}_m(f,x)$ and $\operatorname{orb}_m(f,y)$
were not disjoint, Lemma~\ref{ll:lem:orbit}(f) (with $p=m>0$) would give $y\in\operatorname{orb}_m(f,x)$; so they are
disjoint. Both are contained in $H$, by Step~1 and Lemma~\ref{ll:lem:orbit}(c), and each has exactly $m$ elements by
(ii) and Lemma~\ref{ll:lem:orbit}(e). Hence
\[\card H\ \ge\ \card{\operatorname{orb}_m(f,x)\cup\operatorname{orb}_m(f,y)}=2m\ \ge\ K .\]
Choose $T\subseteq H$ with $\card T=K$. By Lemma~\ref{ll:lem:hv}(a) and Lemma~\ref{fd:lem:mono}(a), $T$ is in convex
position, and $T\subseteq H\subseteq S$. So \lean{HasConvexNGon} $K$ $S$ holds, contradicting (iii). Hence
$y\in\operatorname{orb}_m(f,x)$.

Steps 3 and 4 give $H=\operatorname{orb}_m(f,x)$.
\end{proof}

\begin{lemma}[invariance of $S\setminus\hv(S)$;
\leanref{ErdosSzekeres.LayerLemma.invariant_sdiff_hullVerts}{LAYERLEAN/LayerLemma.lean:217}]\label{ll:lem:inv-sdiff}
Let $f$ be an affine map of $\R^2$ that is injective, $S$ a finite set and $m\in\N$ with $0<m$. Assume $f(y)\in S$ for
every $y\in S$ and $f^m(y)=y$ for every $y\in S$. Then $f(y)\in S\setminus\hv(S)$ for every $y\in S\setminus\hv(S)$.
\end{lemma}

\begin{proof}
We follow the Lean proof. Let $y\in S\setminus\hv(S)$. Then $f(y)\in S$. Suppose $f(y)\in\hv(S)$. By
Lemma~\ref{ll:lem:map-hv}, $f$ maps the finite set $\hv(S)$ into itself, so Lemma~\ref{ll:lem:orbit}(a) gives
$f^{m-1}(f(y))\in\hv(S)$. As $m\ge1$, $f^{m-1}(f(y))=f^m(y)=y$, so $y\in\hv(S)$, a contradiction. Hence
$f(y)\in S\setminus\hv(S)$.
\end{proof}

\begin{remark}[on Theorem~\ref{ll:thm:core}]\label{ll:rem:core}
(1) The proof uses only two geometric facts: an injective affine map that sends $S$ into itself sends $\hv(S)$ into
itself, hence permutes the finite set $\hv(S)$ (Lemma~\ref{ll:lem:map-hv}), and the hull vertices are in convex
position (Lemma~\ref{ll:lem:hv}(a)). No general
position is assumed, and $f$ need not be a rotation.
(2) If $K=0$, hypothesis (iii) is false by (C6), and the theorem is vacuous; if $\hv(S)=\emptyset$, which by
Lemma~\ref{ll:lem:hv}(d) happens only for $S=\emptyset$, there is no $x$.
(3) Under the hypotheses of the theorem, if $S\ne\emptyset$ then $\hv(S)$ contains some $x$ (Lemma~\ref{ll:lem:hv}(d)),
and the conclusion with Lemma~\ref{ll:lem:orbit}(e) and (ii) gives $\card{\hv(S)}=m$.
\end{remark}

\subsection{Rotations}

\begin{definition}[complex coordinates; \leanref{ErdosSzekeres.LayerLemma.toC}{LAYERLEAN/LayerLemma.lean:231},
\leanref{ErdosSzekeres.LayerLemma.ofC}{LAYERLEAN/LayerLemma.lean:234},
\leanref{ErdosSzekeres.LayerLemma.crot}{LAYERLEAN/LayerLemma.lean:252}]\label{ll:def:complex}
For $p\in\R^2$ let $\operatorname{toC}(p)=p_0+ip_1\in\C$; for $z\in\C$ let $\operatorname{ofC}(z)=(\operatorname{Re}z,
\operatorname{Im}z)\in\R^2$; for $\gamma,w\in\C$ and $p\in\R^2$ let
\[\operatorname{crot}_{\gamma,w}(p)=\operatorname{ofC}\bigl(\gamma+w\,(\operatorname{toC}(p)-\gamma)\bigr)\]
(Lean \texttt{crot c w p}, where $c\in\C$ plays the role of $\gamma$).
\end{definition}

\begin{lemma}[complex coordinates; \leanref{ErdosSzekeres.LayerLemma.toC_re}{LAYERLEAN/LayerLemma.lean:236},
\leanref{ErdosSzekeres.LayerLemma.toC_im}{LAYERLEAN/LayerLemma.lean:237},
\leanref{ErdosSzekeres.LayerLemma.ofC_zero}{LAYERLEAN/LayerLemma.lean:238},
\leanref{ErdosSzekeres.LayerLemma.ofC_one}{LAYERLEAN/LayerLemma.lean:239},
\leanref{ErdosSzekeres.LayerLemma.toC_ofC}{LAYERLEAN/LayerLemma.lean:241},
\leanref{ErdosSzekeres.LayerLemma.ofC_toC}{LAYERLEAN/LayerLemma.lean:244},
\leanref{ErdosSzekeres.LayerLemma.toC_injective}{LAYERLEAN/LayerLemma.lean:248},
\leanref{ErdosSzekeres.LayerLemma.toC_crot}{LAYERLEAN/LayerLemma.lean:254},
\leanref{ErdosSzekeres.LayerLemma.crot_one}{LAYERLEAN/LayerLemma.lean:257},
\leanref{ErdosSzekeres.LayerLemma.crot_comp}{LAYERLEAN/LayerLemma.lean:262},
\leanref{ErdosSzekeres.LayerLemma.crot_iterate}{LAYERLEAN/LayerLemma.lean:267},
\leanref{ErdosSzekeres.LayerLemma.crot_eq_self_iff}{LAYERLEAN/LayerLemma.lean:273},
\leanref{ErdosSzekeres.LayerLemma.crot_injective}{LAYERLEAN/LayerLemma.lean:287}]\label{ll:lem:complex}
For all $p\in\R^2$, $z,\gamma,w,v\in\C$ and $k\in\N$:
\begin{enumerate}[label=(\alph*)]
\item (\lean{toC_re}, \lean{toC_im}, \lean{ofC_zero}, \lean{ofC_one}, \lean{toC_ofC}, \lean{ofC_toC},
  \lean{toC_injective}) $\operatorname{Re}\operatorname{toC}(p)=p_0$, $\operatorname{Im}\operatorname{toC}(p)=p_1$,
  $\operatorname{ofC}(z)_0=\operatorname{Re}z$, $\operatorname{ofC}(z)_1=\operatorname{Im}z$,
  $\operatorname{toC}(\operatorname{ofC}(z))=z$ and $\operatorname{ofC}(\operatorname{toC}(p))=p$; in particular
  $\operatorname{toC}$ is injective;
\item (\lean{toC_crot}) $\operatorname{toC}(\operatorname{crot}_{\gamma,w}(p))=\gamma+w(\operatorname{toC}(p)-\gamma)$;
\item (\lean{crot_one}) $\operatorname{crot}_{\gamma,1}=\mathrm{id}$;
\item (\lean{crot_comp}) $\operatorname{crot}_{\gamma,w}(\operatorname{crot}_{\gamma,v}(p))
  =\operatorname{crot}_{\gamma,wv}(p)$;
\item (\lean{crot_iterate}) $(\operatorname{crot}_{\gamma,w})^k=\operatorname{crot}_{\gamma,w^k}$;
\item (\lean{crot_eq_self_iff}) if $w\ne1$, then $\operatorname{crot}_{\gamma,w}(p)=p$ if and only if
  $\operatorname{toC}(p)=\gamma$;
\item (\lean{crot_injective}) if $w\ne0$, then $\operatorname{crot}_{\gamma,w}$ is injective.
\end{enumerate}
\end{lemma}

\begin{proof}
We follow the Lean proofs. (a) The four coordinate identities are the definitions; they give
$\operatorname{toC}(\operatorname{ofC}(z))=\operatorname{Re}z+i\operatorname{Im}z=z$ and
$\operatorname{ofC}(\operatorname{toC}(p))=(p_0,p_1)=p$, and a map with a left inverse is injective.
(b) follows from $\operatorname{toC}\circ\operatorname{ofC}=\mathrm{id}$.
(c) By (b), $\operatorname{toC}(\operatorname{crot}_{\gamma,1}(p))=\gamma+(\operatorname{toC}(p)-\gamma)
=\operatorname{toC}(p)$; apply injectivity of $\operatorname{toC}$.
(d) By (b) twice, $\operatorname{toC}$ of the left side is $\gamma+w\bigl((\gamma+v(\operatorname{toC}(p)-\gamma))
-\gamma\bigr)=\gamma+wv(\operatorname{toC}(p)-\gamma)$, which is $\operatorname{toC}$ of the right side; apply
injectivity.
(e) Induction on $k$: $(\operatorname{crot}_{\gamma,w})^0=\mathrm{id}=\operatorname{crot}_{\gamma,1}$ by (c), and
$(\operatorname{crot}_{\gamma,w})^{k+1}(p)=\operatorname{crot}_{\gamma,w}(\operatorname{crot}_{\gamma,w^k}(p))
=\operatorname{crot}_{\gamma,w\,w^k}(p)$ by (d).
(f) Write $u=\operatorname{toC}(p)$. By (a) and (b), $\operatorname{crot}_{\gamma,w}(p)=p$ iff $\gamma+w(u-\gamma)=u$,
i.e.\ iff $(w-1)(u-\gamma)=0$, i.e.\ (as $w\ne1$ and $\C$ is a field) iff $u=\gamma$.
(g) If $\operatorname{crot}_{\gamma,w}(p)=\operatorname{crot}_{\gamma,w}(q)$, then by (b)
$w(\operatorname{toC}(p)-\operatorname{toC}(q))=0$, so $\operatorname{toC}(p)=\operatorname{toC}(q)$ as $w\ne0$, and
$p=q$ by (a).
\end{proof}

\begin{definition}[rotation about a point; \leanref{ErdosSzekeres.LayerLemma.rotAbout}{LAYERLEAN/LayerLemma.lean:298}]
\label{ll:def:rot}
For $c\in\R^2$ and $\theta\in\R$ let $\rho_{c,\theta}:\R^2\to\R^2$ (Lean \texttt{rotAbout c}~$\theta$) be
\[\rho_{c,\theta}(p)=\bigl(c_0+\cos\theta\,(p_0-c_0)-\sin\theta\,(p_1-c_1),\ \
c_1+\sin\theta\,(p_0-c_0)+\cos\theta\,(p_1-c_1)\bigr).\]
This is the map of Definition~\ref{fd:def:rot}, the counterclockwise rotation by $\theta$ about $c$. For $m\in\N$ the
map $\rho_{c,2\pi/m}$ is Lean's \texttt{rotAbout c (2 * Real.pi / m)}, with the convention (C7) on the angle.
\end{definition}

\begin{lemma}[rotations; \leanref{ErdosSzekeres.LayerLemma.rotAbout_eq_crot}{LAYERLEAN/LayerLemma.lean:302},
\leanref{ErdosSzekeres.LayerLemma.rotAff}{LAYERLEAN/LayerLemma.lean:311},
\leanref{ErdosSzekeres.LayerLemma.rotAbout_injective}{LAYERLEAN/LayerLemma.lean:319},
\leanref{ErdosSzekeres.LayerLemma.rotAbout_iterate}{LAYERLEAN/LayerLemma.lean:465}]\label{ll:lem:rot}
Let $c\in\R^2$, $\theta,\alpha\in\R$ and $j\in\N$.
\begin{enumerate}[label=(\alph*)]
\item (\lean{rotAbout_eq_crot}) $\rho_{c,\theta}=\operatorname{crot}_{\operatorname{toC}(c),\,e^{i\theta}}$.
\item (\lean{rotAff}) For all $p,v\in\R^2$, $\rho_{c,\theta}(v+p)=A_\theta v+\rho_{c,\theta}(p)$, where $A_\theta$ is the
  linear map with matrix $\begin{pmatrix}\cos\theta&-\sin\theta\\ \sin\theta&\cos\theta\end{pmatrix}$ acting on column
  vectors. So $\rho_{c,\theta}$ is an affine map with linear part $A_\theta$; Lean's \texttt{rotAff c}~$\theta$ is this
  affine map, and its underlying map is $\rho_{c,\theta}$.
\item (\lean{rotAbout_injective}) $\rho_{c,\theta}$ is injective.
\item (\lean{rotAbout_iterate}) $(\rho_{c,\alpha})^j=\rho_{c,j\alpha}$ (with $j$ cast to $\R$).
\end{enumerate}
\end{lemma}

\begin{proof}
We follow the Lean proofs.
(a) Let $u=\operatorname{toC}(p)=p_0+ip_1$, $\gamma=\operatorname{toC}(c)=c_0+ic_1$ and
$e^{i\theta}=\cos\theta+i\sin\theta$. Then
\[\gamma+e^{i\theta}(u-\gamma)=\bigl(c_0+\cos\theta(p_0-c_0)-\sin\theta(p_1-c_1)\bigr)
+i\bigl(c_1+\sin\theta(p_0-c_0)+\cos\theta(p_1-c_1)\bigr),\]
whose real and imaginary parts are the coordinates of $\rho_{c,\theta}(p)$. Applying $\operatorname{ofC}$
(Lemma~\ref{ll:lem:complex}(a)) gives $\operatorname{crot}_{\gamma,e^{i\theta}}(p)=\rho_{c,\theta}(p)$.
(b) The formula of Definition~\ref{ll:def:rot} reads $\rho_{c,\theta}(p)=A_\theta(p-c)+c$. Hence
$\rho_{c,\theta}(v+p)=A_\theta(v+p-c)+c=A_\theta v+\rho_{c,\theta}(p)$ by linearity of $A_\theta$.
(c) By (a), $\rho_{c,\theta}=\operatorname{crot}_{\gamma,e^{i\theta}}$ with $e^{i\theta}\ne0$; apply
Lemma~\ref{ll:lem:complex}(g).
(d) By (a) and Lemma~\ref{ll:lem:complex}(e), $(\rho_{c,\alpha})^j=\operatorname{crot}_{\gamma,(e^{i\alpha})^j}$, and
$(e^{i\alpha})^j=e^{ij\alpha}$; by (a) again this is $\rho_{c,j\alpha}$.
\end{proof}

\begin{lemma}[invariance of $S\setminus\{c\}$;
\leanref{ErdosSzekeres.LayerLemma.invariant_erase_centre}{LAYERLEAN/LayerLemma.lean:325}]\label{ll:lem:erase}
Let $S$ be a finite set, $c$ a point and $\theta\in\R$. If $S$ is $\rho_{c,\theta}$-invariant, then $S\setminus\{c\}$
is $\rho_{c,\theta}$-invariant.
\end{lemma}

\begin{proof}
We follow the Lean proof. Write $\rho=\rho_{c,\theta}$. Putting $p=c$ in Definition~\ref{ll:def:rot} gives $\rho(c)=c$.
Let $p\in S\setminus\{c\}$. Then $\rho(p)\in S$. If $\rho(p)=c=\rho(c)$, then $p=c$ by Lemma~\ref{ll:lem:rot}(c), which
is false. Hence $\rho(p)\in S\setminus\{c\}$.
\end{proof}

\begin{lemma}[exact period of the rotation;
\leanref{ErdosSzekeres.LayerLemma.minimalPeriod_rotAbout}{LAYERLEAN/LayerLemma.lean:336}]\label{ll:lem:minperiod}
Let $c$ be a point, $m\in\N$ with $0<m$, and $y$ a point with $y\ne c$. Then
$\operatorname{per}_{\rho}(y)=m$ for $\rho=\rho_{c,2\pi/m}$.
\end{lemma}

\begin{proof}
We follow the Lean proof. Let $\gamma=\operatorname{toC}(c)$ and $\omega=e^{2\pi i/m}$, where $2\pi/m$ is the real
number of (C7). By Lemma~\ref{ll:lem:rot}(a), $\rho=\operatorname{crot}_{\gamma,\omega}$, and by
Lemma~\ref{ll:lem:complex}(e), $\rho^k=\operatorname{crot}_{\gamma,\omega^k}$ for $k\in\N$. We have
$\omega^k=e^{2\pi ik/m}=\cos(2\pi k/m)+i\sin(2\pi k/m)$. Hence $\omega^m=\cos2\pi+i\sin2\pi=1$, and for $0<k<m$ we have
$0<2\pi k/m<2\pi$, so $\cos(2\pi k/m)<1$ and $\omega^k\ne1$. (That is, $\omega$ is a primitive $m$-th root of unity;
Lean uses Mathlib's \lean{Complex.isPrimitiveRoot_exp}.)

By Lemma~\ref{ll:lem:complex}(e) and~(c),
$\rho^m=\operatorname{crot}_{\gamma,\omega^m}=\operatorname{crot}_{\gamma,1}=\mathrm{id}$, so $\rho^m(y)=y$ with $m\ge1$, and
(P2) gives $1\le q\le m$ for $q=\operatorname{per}_\rho(y)$. Suppose $q<m$. By (P1), $\operatorname{crot}_{\gamma,
\omega^q}(y)=\rho^q(y)=y$, where $\omega^q\ne1$ because $0<q<m$. Lemma~\ref{ll:lem:complex}(f) gives
$\operatorname{toC}(y)=\gamma=\operatorname{toC}(c)$, hence $y=c$ by Lemma~\ref{ll:lem:complex}(a), a contradiction.
So $q=m$.
\end{proof}

\begin{theorem}[layer lemma, rotations;
\leanref{ErdosSzekeres.LayerLemma.hullVerts_eq_orbitFS_rotAbout}{LAYERLEAN/LayerLemma.lean:364}]\label{ll:thm:core-rot}
Let $m,K\in\N$ with $0<m$ and $K\le 2m$, let $S$ be a finite set and $c$ a point with $c\notin S$. Assume that $S$ is
$\rho$-invariant for $\rho=\rho_{c,2\pi/m}$ and that $S$ is $K$-free. Then for every $x\in\hv(S)$,
\[\hv(S)=\operatorname{orb}_m(\rho,x)=\{\rho^i(x):\ 0\le i<m\}.\]
\end{theorem}

\begin{proof}
As in Lean, apply Theorem~\ref{ll:thm:core} to the affine map $f=\rho$ (Lemma~\ref{ll:lem:rot}(b); Lean's
\texttt{rotAff c (2 * Real.pi / m)}, whose underlying map is $\rho$), which is injective by Lemma~\ref{ll:lem:rot}(c).
Its hypotheses hold: $0<m$ and $K\le2m$; (i) is the invariance of $S$; for (ii), every $y\in S$ satisfies $y\ne c$
because $c\notin S$, so $\operatorname{per}_\rho(y)=m$ by Lemma~\ref{ll:lem:minperiod}; (iii) is the $K$-freeness of
$S$.
\end{proof}

\subsection{Peeling the layers}

\begin{definition}[peeling; \leanref{ErdosSzekeres.LayerLemma.peel}{LAYERLEAN/LayerLemma.lean:374}]\label{ll:def:peel}
For a finite set $S$ define finite sets $\peel(S,n)$, $n\in\N$, by
\[\peel(S,0)=S,\qquad \peel(S,n+1)=\peel(S,n)\setminus\hv(\peel(S,n)).\]
So $\peel(S,n)$ is $S$ with its $n$ outermost convex layers removed, and $\hv(\peel(S,n))$ is the $n$-th convex layer
of $S$ (counting from $0$).
\end{definition}

\begin{lemma}[peeling, basic facts; \leanref{ErdosSzekeres.LayerLemma.peel_succ}{LAYERLEAN/LayerLemma.lean:378},
\leanref{ErdosSzekeres.LayerLemma.peel_subset}{LAYERLEAN/LayerLemma.lean:381}]\label{ll:lem:peel-basic}
Let $S$ be a finite set. (a) $\peel(S,n+1)=\peel(S,n)\setminus\hv(\peel(S,n))$ for every $n\in\N$. (b)
$\peel(S,n)\subseteq S$ for every $n\in\N$.
\end{lemma}

\begin{proof}
As in Lean: (a) is the definition. (b) By induction: $\peel(S,0)=S$, and
$\peel(S,n+1)\subseteq\peel(S,n)\subseteq S$.
\end{proof}

\begin{lemma}[peeling is antitone; \leanref{ErdosSzekeres.LayerLemma.peel_antitone}{LAYERLEAN/LayerLemma.lean:460}]
\label{ll:lem:antitone}
Let $S$ be a finite set and $n,n'\in\N$ with $n\le n'$. Then $\peel(S,n')\subseteq\peel(S,n)$.
\end{lemma}

\begin{proof}
Induction on $n'\ge n$, as in Lean: for $n'=n$ there is nothing to prove, and
$\peel(S,n'+1)\subseteq\peel(S,n')\subseteq\peel(S,n)$ by Lemma~\ref{ll:lem:peel-basic}(a).
\end{proof}

\begin{lemma}[the peels are invariant; \leanref{ErdosSzekeres.LayerLemma.invariant_peel}{LAYERLEAN/LayerLemma.lean:386}]
\label{ll:lem:inv-peel}
Let $m\in\N$ with $0<m$, let $S$ be a finite set and $c$ a point with $c\notin S$, and assume that $S$ is
$\rho$-invariant for $\rho=\rho_{c,2\pi/m}$. Then $\peel(S,n)$ is $\rho$-invariant for every $n\in\N$.
\end{lemma}

\begin{proof}
Induction on $n$, as in Lean. For $n=0$ this is the hypothesis. Assume $\peel(S,n)$ is $\rho$-invariant. Every
$y\in\peel(S,n)$ lies in $S$ (Lemma~\ref{ll:lem:peel-basic}(b)), so $y\ne c$, so $\operatorname{per}_\rho(y)=m$ by
Lemma~\ref{ll:lem:minperiod} and $\rho^m(y)=y$ by (P1). Lemma~\ref{ll:lem:inv-sdiff}, applied to the finite set
$\peel(S,n)$ and the injective affine map $\rho$ (Lemma~\ref{ll:lem:rot}(b),(c)), shows that $\rho$ maps
$\peel(S,n)\setminus\hv(\peel(S,n))=\peel(S,n+1)$ into itself.
\end{proof}

\begin{theorem}[every layer is one orbit;
\leanref{ErdosSzekeres.LayerLemma.layers_are_orbits}{LAYERLEAN/LayerLemma.lean:402}]\label{ll:thm:layers}
Let $m,K\in\N$ with $0<m$ and $K\le2m$, let $S$ be a finite set and $c$ a point with $c\notin S$, and write
$\rho=\rho_{c,2\pi/m}$. Assume that $S$ is $\rho$-invariant and $K$-free, and let $n\in\N$ be such that $\peel(S,n)$
is nonempty. Then there is $x\in\peel(S,n)$ such that
\begin{enumerate}[label=(\alph*)]
\item $\hv(\peel(S,n))=\operatorname{orb}_m(\rho,x)$;
\item $\card{\operatorname{orb}_m(\rho,x)}=m$;
\item $\peel(S,n+1)\subseteq\conv\bigl(\hv(\peel(S,n))\bigr)$.
\end{enumerate}
\end{theorem}

\begin{proof}
We follow the Lean proof. By Lemma~\ref{ll:lem:hv}(d) there is $x\in\hv(\peel(S,n))$, and $x\in\peel(S,n)$ by
Lemma~\ref{ll:lem:hv-basic}(b). The finite set $U=\peel(S,n)$ satisfies the hypotheses of
Theorem~\ref{ll:thm:core-rot}: $c\notin U$ since $U\subseteq S$ (Lemma~\ref{ll:lem:peel-basic}(b)); $U$ is
$\rho$-invariant (Lemma~\ref{ll:lem:inv-peel}); and $U$ is $K$-free, since a $K$-element subset of $U$ in convex position
would be one of $S$ (Lemma~\ref{fd:lem:mono}(c)). Theorem~\ref{ll:thm:core-rot} gives (a). Since $x\in U$, $x\ne c$, so
$\operatorname{per}_\rho(x)=m$ (Lemma~\ref{ll:lem:minperiod}) and (b) follows from Lemma~\ref{ll:lem:orbit}(e). For (c),
$\peel(S,n+1)=U\setminus\hv(U)\subseteq U\subseteq\conv(\hv(U))$ by Lemma~\ref{ll:lem:hv}(c).
\end{proof}

\begin{lemma}[the layers exhaust $S$; \leanref{ErdosSzekeres.LayerLemma.exists_layer}{LAYERLEAN/LayerLemma.lean:419}]
\label{ll:lem:exists-layer}
Let $S$ be a finite set of points and $y\in S$. Then there is $n\in\N$ with $y\in\hv(\peel(S,n))$.
\end{lemma}

\begin{proof}
We follow the Lean proof. There are no other hypotheses: no symmetry, no general position, no $K$-freeness.

Claim: for every $n$, either $\card{\peel(S,n)}+n\le\card S$ or $\peel(S,n)=\emptyset$. For $n=0$ the first alternative
holds. Assume the claim for $n$. If $\peel(S,n)=\emptyset$, then $\peel(S,n+1)=\emptyset\setminus\hv(\emptyset)
=\emptyset$. Otherwise the first alternative holds for $n$; moreover $\hv(\peel(S,n))$ is a nonempty subset of
$\peel(S,n)$ (Lemmas~\ref{ll:lem:hv}(d) and~\ref{ll:lem:hv-basic}(b)), so
$\card{\peel(S,n+1)}\le\card{\peel(S,n)}-1$ and $\card{\peel(S,n+1)}+(n+1)\le\card{\peel(S,n)}+n\le\card S$.

Taking $n=\card S+1$: if $y\in\peel(S,\card S+1)$, this set is nonempty, so
$\card{\peel(S,\card S+1)}+\card S+1\le\card S$, which is impossible. Hence $y\notin\peel(S,k)$ for some $k$; let $k$ be
the least such index. Since $\peel(S,0)=S\ni y$, $k\ne0$; write $k=n+1$. By minimality $y\in\peel(S,n)$, and
$y\notin\peel(S,n+1)=\peel(S,n)\setminus\hv(\peel(S,n))$, so $y\in\hv(\peel(S,n))$.
\end{proof}

\begin{theorem}[each peel removes exactly $m$ points;
\leanref{ErdosSzekeres.LayerLemma.card_peel}{LAYERLEAN/LayerLemma.lean:1037}]\label{ll:thm:card-peel}
Let $m,K\in\N$ with $0<m$ and $K\le2m$, let $S$ be a finite set and $c$ a point with $c\notin S$, and assume that $S$ is
$\rho_{c,2\pi/m}$-invariant and $K$-free. Let $t\in\N$ with $m\cdot(t\mathbin{\dot{-}}1)<\card S$. Then for every
$n\le t$,
\[\card{\peel(S,n)}+m\cdot n=\card S .\]
\end{theorem}

\begin{proof}
Induction on $n$, as in Lean. For $n=0$: $\card{\peel(S,0)}+0=\card S$. Let $n+1\le t$ and assume
$\card{\peel(S,n)}+mn=\card S$ (the case $n\le t$ of the induction hypothesis). Then $t\ge1$ and
$n\le t-1=t\mathbin{\dot{-}}1$, so $mn\le m(t\mathbin{\dot{-}}1)<\card S$ and $\card{\peel(S,n)}=\card S-mn>0$: the set
$\peel(S,n)$ is nonempty. Theorem~\ref{ll:thm:layers} gives $x$ with $\hv(\peel(S,n))=\operatorname{orb}_m(\rho,x)$ and
$\card{\operatorname{orb}_m(\rho,x)}=m$, so $\card{\hv(\peel(S,n))}=m$. As $\hv(\peel(S,n))\subseteq\peel(S,n)$, we have
$\card{\peel(S,n)}\ge m$ and
\[\card{\peel(S,n+1)}=\card{\peel(S,n)\setminus\hv(\peel(S,n))}=\card{\peel(S,n)}-m .\]
Hence $\card{\peel(S,n+1)}+m(n+1)=\card{\peel(S,n)}-m+mn+m=\card S$.
\end{proof}

The next corollary is not a Lean declaration; it assembles the preceding results into the layer decomposition of a set
of exactly $mt$ points.

\begin{corollary}[layer decomposition; written step]\label{ll:cor:decomp}
Let $m,K,t\in\N$ with $0<m$ and $K\le2m$, let $S$ be a finite set and $c$ a point with $c\notin S$, and write
$\rho=\rho_{c,2\pi/m}$. Assume that $S$ is $\rho$-invariant and $K$-free and that $\card S=mt$. For $n<t$ put
$L_n=\hv(\peel(S,n))$. Then:
\begin{enumerate}[label=(\alph*)]
\item for every $n<t$ there is $x_n\in\peel(S,n)$ with $L_n=\operatorname{orb}_m(\rho,x_n)$ and $\card{L_n}=m$;
\item for every $n\le t$, $\peel(S,n)$ is the disjoint union of $L_n,L_{n+1},\dots,L_{t-1}$; in particular
  $\peel(S,t)=\emptyset$ and $S=L_0\sqcup\dots\sqcup L_{t-1}$;
\item $L_j\subseteq\peel(S,i+1)\subseteq\conv(L_i)$ for all $i<j<t$.
\end{enumerate}
\end{corollary}

\begin{proof}
If $t=0$, then $S=\emptyset$; (a) and (c) are vacuous and (b) says $\peel(S,0)=\emptyset$. Let $t\ge1$. Then
$m(t\mathbin{\dot{-}}1)=mt-m<mt=\card S$, so Theorem~\ref{ll:thm:card-peel} gives $\card{\peel(S,n)}=m(t-n)$ for $n\le t$.
Thus $\peel(S,n)\ne\emptyset$ for $n<t$, and $\peel(S,t)=\emptyset$. (a) is Theorem~\ref{ll:thm:layers}(a),(b). For
(b): for $n<t$, $L_n\subseteq\peel(S,n)$ and $\peel(S,n+1)=\peel(S,n)\setminus L_n$, so $\peel(S,n)$ is the disjoint
union of $L_n$ and $\peel(S,n+1)$; downward induction on $n$, starting from $\peel(S,t)=\emptyset$, gives (b). For (c):
$L_j\subseteq\peel(S,j)\subseteq\peel(S,i+1)$ by Lemma~\ref{ll:lem:antitone} ($i+1\le j$), and
$\peel(S,i+1)\subseteq\conv(L_i)$ by Theorem~\ref{ll:thm:layers}(c), applicable since $\peel(S,i)\ne\emptyset$.
\end{proof}

\subsection{Reduction to exactly \texorpdfstring{$mt$}{mt} points and the octagon theorems}

\begin{theorem}[reduction to exactly $mt$ points;
\leanref{ErdosSzekeres.LayerLemma.convexNGon_of_layers}{LAYERLEAN/LayerLemma.lean:1066}]\label{ll:thm:reduction}
Let $m,K,t\in\N$ with $0<m$ and $K\le 2m$, and assume the premise
\begin{itemize}
\item[(H)] for every finite set $T$ and every point $c'$: if $T$ is in general position, $c'\notin T$, $T$ is
  $\rho_{c',2\pi/m}$-invariant and $\card T=m\cdot t$, then \lean{HasConvexNGon} $K$ $T$.
\end{itemize}
Then for every finite set $S$ and every point $c$: if $S$ is in general position, $S$ is $\rho_{c,2\pi/m}$-invariant,
and $m\cdot(t\mathbin{\dot{-}}1)<\card{S\setminus\{c\}}$, then \lean{HasConvexNGon} $K$ $S$.
\end{theorem}

In Lean the bound centre in (H) is also called $c$; we rename it $c'$. The size condition is on $S\setminus\{c\}$
(Lean \texttt{S.erase c}), not on $S$, and $c\in S$ is allowed. The premise (H) is an explicit hypothesis of the
theorem, and no theorem of \dirn{LAYERLEAN/LayerLemma.lean} proves it. In the two applications below it is supplied by
the propositions \lean{Sixfold66} and \lean{Sevenfold70}, which are themselves hypotheses (Remark~\ref{ll:rem:hyp}).

\begin{proof}
We follow the Lean proof. Write $\rho=\rho_{c,2\pi/m}$ and suppose that $S$ is $K$-free. Put $S'=S\setminus\{c\}$. Then
$c\notin S'$; $S'$ is $\rho$-invariant by Lemma~\ref{ll:lem:erase}; and $S'$ is $K$-free by
Lemma~\ref{fd:lem:mono}(c), since $S'\subseteq S$.

Theorem~\ref{ll:thm:card-peel} applies to $S'$ with the given $t$ (its size hypothesis $m(t\mathbin{\dot{-}}1)<\card{S'}$
is our hypothesis on $\card{S\setminus\{c\}}$), and with $n=t$ it gives $\card{\peel(S',t)}+mt=\card{S'}$.
Let $T=S'\setminus\peel(S',t)$. Since $\peel(S',t)\subseteq S'$ (Lemma~\ref{ll:lem:peel-basic}(b)),
$\card T=\card{S'}-\card{\peel(S',t)}=mt$.

$T$ is $\rho$-invariant. Let $y\in T$, so $y\in S'$ and $y\notin\peel(S',t)$. Then $\rho(y)\in S'$. Suppose
$\rho(y)\in\peel(S',t)$. By Lemma~\ref{ll:lem:inv-peel} (applied to $S'$; $c\notin S'$) the set $\peel(S',t)$ is
$\rho$-invariant, so $\rho^{m-1}(\rho(y))\in\peel(S',t)$ by Lemma~\ref{ll:lem:orbit}(a). As $y\in S'$, $y\ne c$, so
$\operatorname{per}_\rho(y)=m$ (Lemma~\ref{ll:lem:minperiod}) and, by (P1) and $m\ge1$,
$\rho^{m-1}(\rho(y))=\rho^m(y)=y$. Hence $y\in\peel(S',t)$, a contradiction. So $\rho(y)\in T$.

Moreover $T\subseteq S'\subseteq S$, so $T$ is in general position (Lemma~\ref{fd:lem:mono}(b)), and $c\notin T$. By (H)
with $c'=c$, \lean{HasConvexNGon} $K$ $T$, hence \lean{HasConvexNGon} $K$ $S$ (Lemma~\ref{fd:lem:mono}(c)),
contradicting the assumption. Hence \lean{HasConvexNGon} $K$ $S$.
\end{proof}

\begin{definition}[the propositions \lean{Sixfold66} and \lean{Sevenfold70};
\leanref{ErdosSzekeres.LayerLemma.Sixfold66}{LAYERLEAN/LayerLemma.lean:1102},
\leanref{ErdosSzekeres.LayerLemma.Sevenfold70}{LAYERLEAN/LayerLemma.lean:1107}]\label{ll:def:66-70}
\lean{Sixfold66} is the proposition: for every finite set $T$ and every point $c$, if $T$ is in general position,
$c\notin T$, $T$ is $\rho_{c,2\pi/6}$-invariant and $\card T=66$, then \lean{HasConvexNGon} $8$ $T$.

\lean{Sevenfold70} is the proposition: for every finite set $T$ and every point $c$, if $T$ is in general position,
$c\notin T$, $T$ is $\rho_{c,2\pi/7}$-invariant and $\card T=70$, then \lean{HasConvexNGon} $8$ $T$.

In both, the angle is written with a real numeral, \texttt{2 * Real.pi / 6} and \texttt{2 * Real.pi / 7}
(Remark~\ref{ll:rem:numerals}). Both are definitions of propositions; neither is proved in Lean
(Remark~\ref{ll:rem:hyp}).
\end{definition}

\begin{theorem}[sixfold octagon theorem, conditional on \lean{Sixfold66};
\leanref{ErdosSzekeres.LayerLemma.sixfold_octagon}{LAYERLEAN/LayerLemma.lean:1113}]\label{ll:thm:sixfold}
Assume the proposition \lean{Sixfold66} (a hypothesis, not proved in Lean; Remark~\ref{ll:rem:hyp}). Let $S$ be a finite
set of points and $c$ a point such that $S$ is in general position, $S$ is $\rho_{c,2\pi/6}$-invariant, and
$61\le\card S$. Then \lean{HasConvexNGon} $8$ $S$: some $8$ points of $S$ are in convex position.
\end{theorem}

\begin{proof}
We follow the Lean proof. Write $\rho=\rho_{c,2\pi/6}$; by Remark~\ref{ll:rem:numerals} this is $\rho_{c,2\pi/m}$ for the
natural number $m=6$.

Step 1. There is $p\in S$ with $p\ne c$: otherwise $S\subseteq\{c\}$ and $\card S\le1<61$.

Step 2. $c\notin S$. This is Lemma~\ref{lg:lem:centre} (Lean \lean{centre_not_mem_of_even}) with $m=6$, which is
positive and even, applied to the $\rho$-invariant set $S$ in general position and the point $p$ of Step~1. For
completeness, the argument in this case: suppose $c\in S$. By Lemma~\ref{ll:lem:rot}(d),
$\rho^3=\rho_{c,3\cdot2\pi/6}=\rho_{c,\pi}$, and since $\cos\pi=-1$, $\sin\pi=0$, Definition~\ref{ll:def:rot} gives
$q:=\rho^3(p)=2c-p$, which lies in $S$ by Lemma~\ref{ll:lem:orbit}(a). The points $p,c,q$ of $S$ are pairwise distinct:
$p\ne c$; $q=c$ and $q=p$ would each give $p=c$. But, by the formula of Definition~\ref{fd:def:orient},
\[\orient(p,c,2c-p)=(c_0-p_0)\cdot2(c_1-p_1)-(c_1-p_1)\cdot2(c_0-p_0)=0,\]
so $\{p,c,q\}$ is collinear by Lemma~\ref{fd:lem:collinear}, contradicting general position
(Definition~\ref{fd:def:gp}).

Step 3. Apply Theorem~\ref{ll:thm:reduction} with $m=6$, $K=8$, $t=11$: $0<6$ and $8\le12$. Premise (H) holds: for $T$
and $c'$ as there (the rotation $\rho_{c',2\pi/m}$ with $m=6$ is $\rho_{c',2\pi/6}$, Remark~\ref{ll:rem:numerals}),
$\card T=6\cdot11=66$, and \lean{Sixfold66} gives \lean{HasConvexNGon} $8$ $T$. $S$ is in general
position and $\rho$-invariant. By Step~2, $S\setminus\{c\}=S$, so $6\cdot(11\mathbin{\dot{-}}1)=60<61\le\card S
=\card{S\setminus\{c\}}$. Theorem~\ref{ll:thm:reduction} gives \lean{HasConvexNGon} $8$ $S$.
\end{proof}

\begin{theorem}[sevenfold octagon theorem, conditional on \lean{Sevenfold70};
\leanref{ErdosSzekeres.LayerLemma.sevenfold_octagon}{LAYERLEAN/LayerLemma.lean:1136}]\label{ll:thm:sevenfold}
Assume the proposition \lean{Sevenfold70} (a hypothesis, not proved in Lean; Remark~\ref{ll:rem:hyp}). Let $S$ be a
finite set of points and $c$ a point such that $S$ is in general position, $S$ is $\rho_{c,2\pi/7}$-invariant, and
$65\le\card S$. Then \lean{HasConvexNGon} $8$ $S$. (The centre $c$ may belong to $S$.)
\end{theorem}

\begin{proof}
We follow the Lean proof. Apply Theorem~\ref{ll:thm:reduction} with $m=7$, $K=8$, $t=10$ (Remark~\ref{ll:rem:numerals}
for the angle): $0<7$ and $8\le14$. Premise (H) holds: for $T$ and $c'$ as there (the rotation is
$\rho_{c',2\pi/7}$), $\card T=7\cdot10=70$, and
\lean{Sevenfold70} gives \lean{HasConvexNGon} $8$ $T$. $S$ is in general position and $\rho_{c,2\pi/7}$-invariant, and
$\card{S\setminus\{c\}}\ge\card S-1\ge64>63=7\cdot(10\mathbin{\dot{-}}1)$. Theorem~\ref{ll:thm:reduction} gives
\lean{HasConvexNGon} $8$ $S$. No exclusion of the centre is needed.
\end{proof}

\begin{remark}[status of \lean{Sixfold66} and \lean{Sevenfold70}]\label{ll:rem:hyp}
\lean{Sixfold66} and \lean{Sevenfold70} are Lean definitions of propositions (\texttt{def ... : Prop}), neither theorems
nor axioms. Nothing in \dirn{LAYERLEAN/LayerLemma.lean} proves them; Theorems~\ref{ll:thm:sixfold}
and~\ref{ll:thm:sevenfold} take them as explicit hypotheses, so they do not enter the axiom lists of these theorems.
Among the Lean files of the repository, the only theorems whose conclusion is \lean{Sixfold66} are in
\dirn{ENCBLEAN/geo/ES8Geo/Main.lean} and take further hypotheses (see below), and no theorem has the conclusion
\lean{Sevenfold70}. What is Lean-checked is the implication: \lean{Sixfold66} implies that every
$\rho_{c,2\pi/6}$-invariant set of at least $61$ points in general position has $8$ points in convex position, and
\lean{Sevenfold70} implies the same for $\rho_{c,2\pi/7}$ and at least $65$ points. These octagon statements are
therefore established here only conditionally.

In the project, the two propositions are supported as follows; none of this is part of the present section. The
record \dirn{C6UPPER} gives the unsatisfiability of SAT encodings of octagon-free configurations of eleven nested
regular hexagons in general position, by cake\_lpr verdicts on cube formulas whose cover is checked there: the formula
of encoder~B (\dirn{C6UPPER/encb/gen_b2.py}) alone, through $64$ cubes, and that of encoder~A
(\dirn{C7UPPER/gen.py}) with an added symmetry-breaking unit clause, through $32$ cubes. The record \dirn{C7UPPER}
gives the unsatisfiability of two separately implemented encodings of the same kind for ten nested regular heptagons
(encoder~A's with an added symmetry-breaking unit clause, encoder~B's alone), each by a cake\_lpr verdict. In the
terms of this document these unsatisfiability results are \emph{computer-assisted}: they rest on LRAT proofs checked
by cake\_lpr, outside Lean. The project's records call such a result a ``certificate'' (an unsatisfiability claim whose
LRAT proof cake\_lpr accepted; \dirn{WRITEUP/note/symmetric_extremal.tex}, lines 109--114); that is a label of the
project, not one of the labels of this document. The step from these unsatisfiability results to \lean{Sixfold66} and
\lean{Sevenfold70} is the written soundness proof of the encodings in \dirn{WRITEUP/note/k8_soundness.tex} (a written
proof, not refereed), which also covers the symmetry breaking and the
cube cover. That note lists the gaps that remain, (G1)--(G4): the correspondence between the generating programs and
their clauses is established by reading the source code; the passage from the formulas to \lean{Sixfold66} and
\lean{Sevenfold70} is a written proof; the angular part of encoder~A rests on binary64 floating-point arithmetic;
and several of its written proofs have had no internal review. For $m=6$ the unit \dirn{ENCBLEAN} also contains a Lean
derivation of \lean{Sixfold66}, \lean{EncB.Geo.sixfold66_of_cubes64}; its \lean{Sixfold66} is that of
\dirn{ENCBLEAN/geo/ES8Geo/Capstone2.lean}, which contains a verbatim copy of the declarations of this file (over a
byte-identical copy of \dirn{ErdosSzekeres/Definitions.lean}). Its hypothesis \lean{EncB.CubesUnsat64} is the
unsatisfiability of the $64$ cube formulas built from a Lean definition of encoder~B's formula (\lean{EncB.encB}).
Neither the identity of the printed texts of these formulas with the cube files that carry the cake\_lpr verdicts nor
the reading of those verdicts as unsatisfiability is Lean-checked; both rest on the trust items (E1)--(E6) listed
under gap~(G2) of that note.

The link with nested polygons is Corollary~\ref{ll:cor:decomp}: a set $T$ that violates \lean{Sixfold66} is in general
position, centre-free, $\rho_{c,2\pi/6}$-invariant, $8$-free and has $66=6\cdot11$ points, so (with $m=6$, $K=8\le12$,
$t=11$) it is the disjoint union of eleven $\rho_{c,2\pi/6}$-orbits $L_0,\dots,L_{10}$ of $6$ points each, with
$L_j\subseteq\conv(L_i)$ for $i<j$; likewise a set violating \lean{Sevenfold70} is the disjoint union of ten orbits of
$7$ points.
\end{remark}

\begin{remark}[real-cast numerals]\label{ll:rem:numerals}
In the general statements of this section the angle is \texttt{2 * Real.pi / m} with $m\in\N$ cast to $\R$ (C7). In
\lean{Sixfold66}, \lean{Sevenfold70}, \lean{sixfold_octagon} and \lean{sevenfold_octagon} the angle is
\texttt{2 * Real.pi / 6}, resp.\ \texttt{2 * Real.pi / 7}, where $6$ and $7$ are real numerals. The cast of the
natural number
$6$ is the real number $6$ (likewise for $7$), so $\rho_{c,2\pi/6}$ is the same map in both readings; the Lean proofs
bridge the two forms with \texttt{exact\_mod\_cast}, and the cardinality conditions $\card T=6\cdot 11$ and $\card T=66$
(resp.\ $7\cdot10$ and $70$) with \texttt{simpa}. Mathematically nothing changes.
\end{remark}

\paragraph{Coverage.}
All declarations of this section's scope in \dirn{LAYERLEAN/LayerLemma.lean} are covered:
\lean{orbitFS} (Definition~\ref{ll:def:orbit});
\lean{iterate_mem}, \lean{mem_orbitFS_self}, \lean{orbitFS_subset}, \lean{orbitFS_invariant},
\lean{card_orbitFS_of_minimalPeriod}, \lean{mem_orbitFS_of_not_disjoint} (Lemma~\ref{ll:lem:orbit});
\lean{Invariant} (Definition~\ref{ll:def:inv});
\lean{hullVerts} (Definition~\ref{ll:def:hv}); \lean{mem_hullVerts}, \lean{hullVerts_subset}
(Lemma~\ref{ll:lem:hv-basic}); \lean{convexPosition_hullVerts}, \lean{extreme_mem_hullVerts},
\lean{subset_convexHull_hullVerts}, \lean{hullVerts_nonempty} (Lemma~\ref{ll:lem:hv}; the proofs of (b)--(d) are
elementary and differ from the Lean proofs, which use two Mathlib facts on extreme points for (b), the Krein--Milman
theorem for (c) and the Krein--Milman lemma for (d));
\lean{map_mem_hullVerts} (Lemma~\ref{ll:lem:map-hv});
\lean{hullVerts_eq_orbitFS} (Theorem~\ref{ll:thm:core});
\lean{invariant_sdiff_hullVerts} (Lemma~\ref{ll:lem:inv-sdiff});
\lean{toC}, \lean{ofC}, \lean{crot} (Definition~\ref{ll:def:complex});
\lean{toC_re}, \lean{toC_im}, \lean{ofC_zero}, \lean{ofC_one}, \lean{toC_ofC}, \lean{ofC_toC}, \lean{toC_injective},
\lean{toC_crot}, \lean{crot_one}, \lean{crot_comp}, \lean{crot_iterate}, \lean{crot_eq_self_iff}, \lean{crot_injective}
(Lemma~\ref{ll:lem:complex});
\lean{rotAbout} (Definition~\ref{ll:def:rot});
\lean{rotAbout_eq_crot}, \lean{rotAff}, \lean{rotAbout_injective}, \lean{rotAbout_iterate} (Lemma~\ref{ll:lem:rot});
\lean{invariant_erase_centre} (Lemma~\ref{ll:lem:erase});
\lean{minimalPeriod_rotAbout} (Lemma~\ref{ll:lem:minperiod});
\lean{hullVerts_eq_orbitFS_rotAbout} (Theorem~\ref{ll:thm:core-rot});
\lean{peel} (Definition~\ref{ll:def:peel}); \lean{peel_succ}, \lean{peel_subset} (Lemma~\ref{ll:lem:peel-basic});
\lean{peel_antitone} (Lemma~\ref{ll:lem:antitone});
\lean{invariant_peel} (Lemma~\ref{ll:lem:inv-peel});
\lean{layers_are_orbits} (Theorem~\ref{ll:thm:layers});
\lean{exists_layer} (Lemma~\ref{ll:lem:exists-layer});
\lean{card_peel} (Theorem~\ref{ll:thm:card-peel});
\lean{convexNGon_of_layers} (Theorem~\ref{ll:thm:reduction});
\lean{Sixfold66}, \lean{Sevenfold70} (Definition~\ref{ll:def:66-70});
\lean{sixfold_octagon} (Theorem~\ref{ll:thm:sixfold}, conditional on \lean{Sixfold66});
\lean{sevenfold_octagon} (Theorem~\ref{ll:thm:sevenfold}, conditional on \lean{Sevenfold70}).
Corollary~\ref{ll:cor:decomp} is a written step, not a Lean declaration. Nothing in the scope is deliberately left out.
The declarations of lines 471--1032 of the same file (from \lean{orient_rotAbout} to \lean{dist2_le_of_later}) are
outside this section's scope; \lean{rotAbout_iterate} (line 465) is included here because it lies in the range of this
section, and Step~2 of our proof of Theorem~\ref{ll:thm:sixfold} uses it.
